% !TeX root = ../main.tex

\xchapter{相关理论与技术}{Related Theory and Technology}

\xsection{强化学习与连续控制基础理论}{Fundamental Theory of Reinforcement Learning and Continuous Control}

强化学习（Reinforcement Learning, RL）是机器学习的重要分支，旨在研究智能体（Agent）如何通过与环境（Environment）的序贯交互，依据累积奖励信号逐步习得使长期回报最大化的最优决策策略。本节将系统阐述连续控制场景下强化学习的核心数学框架：首先定义马尔可夫决策过程及其在部分可观测情形下的扩展形式，继而从策略梯度定理出发推导 Actor-Critic 架构的理论依据，最后深入剖析最大熵强化学习框架及 Soft Actor-Critic 算法的核心机制，为后续章节的算法设计与表征优化提供理论基石。

\xsubsection{马尔可夫决策过程与部分可观测马尔可夫决策过程}{Markov Decision Process and Partially Observable Markov Decision Process}

\textbf{马尔可夫决策过程}

强化学习问题的标准数学抽象为马尔可夫决策过程（Markov Decision Process, MDP）。一个完整的 MDP 由五元组 $\mathcal{M} = \langle \mathcal{S}, \mathcal{A}, \mathcal{P}, \mathcal{R}, \gamma \rangle$ 严格定义，各分量含义如下：
\begin{itemize}
    \item $\mathcal{S}$ 为状态空间（State Space），表示环境中所有可能状态的集合。在连续控制任务中，$\mathcal{S} \subseteq \mathbb{R}^n$ 通常为高维连续空间，描述系统的完整物理配置（如关节角度、速度、接触力等）。
    \item $\mathcal{A}$ 为动作空间（Action Space），表示智能体可执行动作的集合。连续控制任务中，$\mathcal{A} \subseteq \mathbb{R}^m$ 为连续向量空间，对应电机力矩或关节速度等控制指令。
    \item $\mathcal{P}: \mathcal{S} \times \mathcal{A} \times \mathcal{S} \to [0, 1]$ 为状态转移概率函数，$\mathcal{P}(s' \mid s, a)$ 描述智能体在状态 $s \in \mathcal{S}$ 执行动作 $a \in \mathcal{A}$ 后转移至状态 $s' \in \mathcal{S}$ 的概率，刻画了环境的随机动力学特性。
    \item $\mathcal{R}: \mathcal{S} \times \mathcal{A} \to \mathbb{R}$ 为奖励函数，$r(s, a) = \mathcal{R}(s, a)$ 定义了智能体在状态 $s$ 下执行动作 $a$ 所获得的即时标量反馈，是引导策略优化的核心信号。
    \item $\gamma \in [0, 1)$ 为折扣因子（Discount Factor），用于权衡未来奖励与即时奖励的相对重要性，同时确保无限时域下累积回报的数学收敛性。
\end{itemize}

MDP 的核心假设为马尔可夫性质（Markov Property）：系统下一状态的条件概率分布仅依赖于当前状态与动作，而与历史轨迹无关，即
\begin{equation}
    \mathcal{P}(s_{t+1} \mid s_t, a_t, s_{t-1}, a_{t-1}, \ldots, s_0, a_0) = \mathcal{P}(s_{t+1} \mid s_t, a_t)
\end{equation}

在每一时间步 $t$，智能体观测到当前状态 $s_t$，依据策略 $\pi: \mathcal{S} \to \Delta(\mathcal{A})$ 选择动作 $a_t \sim \pi(\cdot \mid s_t)$（其中 $\Delta(\mathcal{A})$ 表示 $\mathcal{A}$ 上的概率单纯形），环境转移至新状态 $s_{t+1} \sim \mathcal{P}(\cdot \mid s_t, a_t)$ 并产生即时奖励 $r_t = r(s_t, a_t)$。智能体的优化目标是寻找最优策略 $\pi^*$，使得折扣累积期望回报最大化：
\begin{equation}
    J(\pi) = \mathbb{E}_{\tau \sim \pi} \left[ \sum_{t=0}^{\infty} \gamma^t r(s_t, a_t) \right]
\end{equation}
其中，$\tau = (s_0, a_0, s_1, a_1, \ldots)$ 为策略 $\pi$ 与环境交互产生的轨迹（Trajectory）。

为评估策略的优劣，定义状态价值函数（State Value Function）$V^\pi: \mathcal{S} \to \mathbb{R}$ 与动作价值函数（Action-Value Function）$Q^\pi: \mathcal{S} \times \mathcal{A} \to \mathbb{R}$：
\begin{equation}
    V^\pi(s) = \mathbb{E}_{\pi} \left[ \sum_{k=0}^{\infty} \gamma^k r_{t+k} \;\middle|\; s_t = s \right]
\end{equation}
\begin{equation}
    Q^\pi(s, a) = \mathbb{E}_{\pi} \left[ \sum_{k=0}^{\infty} \gamma^k r_{t+k} \;\middle|\; s_t = s,\, a_t = a \right]
\end{equation}

两者之间满足 $V^\pi(s) = \mathbb{E}_{a \sim \pi(\cdot|s)}[Q^\pi(s, a)]$。根据动态规划原理，$Q^\pi$ 满足贝尔曼期望方程（Bellman Expectation Equation）：
\begin{equation}
    Q^\pi(s, a) = \mathbb{E}_{s' \sim \mathcal{P}(\cdot|s,a)} \left[ r(s, a) + \gamma \mathbb{E}_{a' \sim \pi(\cdot|s')} \left[ Q^\pi(s', a') \right] \right]
\end{equation}

此递归关系构成了基于时序差分（Temporal Difference, TD）学习的理论基础，同时也是 Critic 网络参数更新的核心依据。

\textbf{部分可观测马尔可夫决策过程}

标准 MDP 假设智能体在每一时间步均可精确获得系统的完整状态 $s_t$。然而，在大量现实控制场景中，智能体仅能获取原始传感器数据（如图像像素、激光雷达点云或局部关节读数），这些数据往往只是真实系统状态的不完整投影，且可能包含噪声。此类场景更为精确的数学描述为部分可观测马尔可夫决策过程（Partially Observable Markov Decision Process, POMDP），由七元组 $\mathcal{M}_{\mathrm{po}} = \langle \mathcal{S}, \mathcal{A}, \mathcal{P}, \mathcal{R}, \Omega, \mathcal{O}, \gamma \rangle$ 定义，其中 $\Omega$ 为观测空间（Observation Space），$\mathcal{O}: \mathcal{S} \times \mathcal{A} \to \Delta(\Omega)$ 为观测概率函数，$\mathcal{O}(o \mid s', a)$ 表示智能体执行动作 $a$ 并转移至状态 $s'$ 后获得观测 $o \in \Omega$ 的概率。

在 POMDP 中，观测 $o_t$ 与状态 $s_t$ 的本质区别在于：$s_t$ 包含完整刻画系统动力学所需的全部信息，而 $o_t$ 仅为 $s_t$ 的一个条件采样，即
\begin{equation}
    o_t \sim \mathcal{O}(\cdot \mid s_t, a_{t-1})
\end{equation}

由于 $o_t$ 无法唯一确定 $s_t$，智能体必须依赖历史交互序列 $h_t = (o_0, a_0, o_1, a_1, \ldots, o_t)$ 来推断状态的后验分布。这一推断在数学上以信念状态（Belief State）$b_t \in \Delta(\mathcal{S})$ 表示：
\begin{equation}
    b_t(s) = \mathcal{P}(s_t = s \mid h_t)
\end{equation}

信念状态依据贝叶斯滤波（Bayesian Filtering）递归更新：
\begin{equation}
    b_{t+1}(s') = \eta \cdot \mathcal{O}(o_{t+1} \mid s', a_t) \sum_{s \in \mathcal{S}} \mathcal{P}(s' \mid s, a_t)\, b_t(s)
\end{equation}
其中，$\eta$ 为归一化常数，确保 $b_{t+1}$ 在状态空间上积分为 1。

以基于图像的视觉控制任务为例，真实状态 $s_t$ 包含机器人的完整关节配置与速度信息，而智能体仅能获取 RGB 图像 $o_t \in \mathbb{R}^{H \times W \times 3}$。单帧图像不仅无法反映速度等动态量，还可能受到遮挡、光照变化与传感器噪声的影响。这一"状态不确定性"（State Uncertainty）使得直接将 $o_t$ 代入为 MDP 中的 $s_t$ 存在根本性的理论缺陷：策略 $\pi(a \mid o_t)$ 所依赖的信息量不足以满足马尔可夫性质，进而导致值函数估计的系统性偏差与策略收敛性能的下降。

上述分析表明，在部分可观测场景下，有效的状态表征（State Representation）是算法性能的关键瓶颈：一方面，需要从高维观测中提取紧凑而充分的特征；另一方面，需要跨时间步聚合历史信息以近似马尔可夫状态。这一内在需求为后文引入自监督表示学习与序列建模技术提供了直接的理论动机。

\xsubsection{Actor-Critic 架构}{Actor-Critic Architecture}

Actor-Critic 架构是连续动作空间控制任务中最为主流的深度强化学习范式之一。其理论基础源于策略梯度定理（Policy Gradient Theorem），通过将策略优化（Actor）与值函数估计（Critic）相解耦，在保留策略梯度方法灵活性的同时，系统性地降低了梯度估计的方差。

\textbf{策略梯度定理}

设参数化随机策略为 $\pi_\theta(a \mid s)$，其中 $\theta$ 为可学习参数。策略的优化目标为最大化期望回报：
\begin{equation}
    J(\theta) = \mathbb{E}_{\tau \sim \pi_\theta} \left[ \sum_{t=0}^{\infty} \gamma^t r(s_t, a_t) \right] = \mathbb{E}_{s \sim d^{\pi_\theta}} \left[ V^{\pi_\theta}(s) \right]
\end{equation}
其中，$d^{\pi_\theta}(s) = \sum_{t=0}^{\infty} \gamma^t \mathcal{P}(s_t = s \mid \pi_\theta)$ 为策略 $\pi_\theta$ 诱导的折扣状态访问分布（Discounted State Visitation Distribution）。

策略梯度定理给出了目标函数关于策略参数的解析梯度形式。利用对数导数技巧（Log-Derivative Trick）展开，并结合贝尔曼方程递归代入，可严格推导得：
\begin{equation}
    \nabla_\theta J(\theta) = \mathbb{E}_{s \sim d^{\pi_\theta},\, a \sim \pi_\theta(\cdot|s)} \left[ \nabla_\theta \log \pi_\theta(a \mid s) \cdot Q^{\pi_\theta}(s, a) \right]
\end{equation}

此公式的核心意义在于：无需对环境动力学 $\mathcal{P}$ 求导，梯度仅通过对策略的微分即可获得，从而使得策略梯度方法在模型未知的情形下依然适用。

\textbf{基线方差缩减与优势函数}

直接使用 $Q^{\pi_\theta}(s, a)$ 作为梯度权重的 REINFORCE 算法，尽管在期望意义下无偏，但 $Q$ 值的绝对量级差异较大，导致梯度估计具有极高方差，使训练不稳定。策略梯度定理允许在不改变梯度期望值的前提下，引入任意与动作 $a$ 无关的基线函数 $b(s)$：
\begin{equation}
    \nabla_\theta J(\theta) = \mathbb{E}_{s \sim d^{\pi_\theta},\, a \sim \pi_\theta(\cdot|s)} \left[ \nabla_\theta \log \pi_\theta(a \mid s) \cdot \left( Q^{\pi_\theta}(s, a) - b(s) \right) \right]
\end{equation}

基线不改变梯度期望的无偏性由下式保证：
\begin{equation}
    \mathbb{E}_{a \sim \pi_\theta(\cdot|s)} \left[ \nabla_\theta \log \pi_\theta(a \mid s) \cdot b(s) \right] = b(s)\, \nabla_\theta \underbrace{\sum_{a} \pi_\theta(a \mid s)}_{=\,1} = 0
\end{equation}

选取 $b(s) = V^{\pi_\theta}(s)$ 可将梯度方差缩减至接近理论下界。此时，梯度权重退化为优势函数（Advantage Function）：
\begin{equation}
    A^{\pi_\theta}(s, a) = Q^{\pi_\theta}(s, a) - V^{\pi_\theta}(s)
\end{equation}

$A^{\pi_\theta}(s, a)$ 的物理含义为：在状态 $s$ 下执行动作 $a$ 相较于当前策略的平均水平所带来的额外收益。当 $A > 0$ 时，该动作优于平均，策略应增大其选择概率；当 $A < 0$ 时，策略应降低其概率，从而实现精确的信用分配（Credit Assignment）。

\textbf{Actor-Critic 的双网络协同优化}

Actor-Critic 架构通过引入两个独立的参数化网络，将策略优化与值函数估计解耦：
\begin{itemize}
    \item \textbf{Actor 网络} $\pi_\theta(a \mid s)$：参数化随机策略，负责生成动作。在连续控制中，通常输出高斯分布的均值 $\mu_\theta(s)$ 与标准差 $\sigma_\theta(s)$，即 $a \sim \mathcal{N}(\mu_\theta(s), \sigma_\theta^2(s)\mathbf{I})$。
    \item \textbf{Critic 网络} $V_\phi(s)$ 或 $Q_\phi(s, a)$：参数化值函数，负责估计当前策略的回报预期，为 Actor 提供低方差的梯度信号。
\end{itemize}

Critic 通过最小化时序差分误差进行参数更新，以单步 TD 为例，其损失函数定义为：
\begin{equation}
    L_{\mathrm{critic}}(\phi) = \mathbb{E}_{(s_t, a_t, r_t, s_{t+1}) \sim \mathcal{D}} \left[ \left( V_\phi(s_t) - \left( r_t + \gamma V_{\bar{\phi}}(s_{t+1}) \right) \right)^2 \right]
\end{equation}
其中，$V_{\bar{\phi}}$ 表示目标网络（Target Network），其参数通过指数移动平均缓慢跟踪 $\phi$，以稳定自举（Bootstrapping）过程。Actor 利用 Critic 估计的优势函数执行策略梯度更新：
\begin{equation}
    \nabla_\theta J(\theta) \approx \mathbb{E}_{(s_t, a_t) \sim \mathcal{D}} \left[ \nabla_\theta \log \pi_\theta(a_t \mid s_t) \cdot \hat{A}(s_t, a_t) \right]
\end{equation}
其中，$\hat{A}(s_t, a_t) = r_t + \gamma V_{\bar{\phi}}(s_{t+1}) - V_\phi(s_t)$ 为 TD 残差对优势函数的单步估计。

Actor-Critic 架构的本质贡献在于将高方差的蒙特卡洛回报估计替换为低方差的自举估计，在样本效率与估计精度之间取得有效平衡。然而，该架构中 Critic 网络的估计质量直接制约着 Actor 梯度信号的可靠性：若 Critic 基于不充分的状态表征进行值函数估计，则梯度偏差将直接损害策略的收敛性能。这一依赖关系深刻揭示了表征质量对整体算法性能的核心影响。

\xsubsection{最大熵强化学习与 Soft Actor-Critic 算法}{Maximum Entropy Reinforcement Learning and Soft Actor-Critic}

\textbf{最大熵强化学习目标}

标准强化学习仅优化期望累积回报，当策略趋于确定性时，智能体的探索能力随之退化，极易陷入局部最优。最大熵强化学习（Maximum Entropy Reinforcement Learning）在优化目标中引入策略熵正则项，将目标函数修正为：
\begin{equation}
    J_{\mathrm{MaxEnt}}(\pi) = \sum_{t=0}^{T} \mathbb{E}_{(s_t, a_t) \sim \rho_\pi} \left[ r(s_t, a_t) + \alpha \mathcal{H}\!\left( \pi(\cdot \mid s_t) \right) \right]
\end{equation}
其中，$\mathcal{H}(\pi(\cdot \mid s_t)) = -\mathbb{E}_{a \sim \pi(\cdot|s_t)}[\log \pi(a \mid s_t)]$ 为策略在状态 $s_t$ 下的香农熵（Shannon Entropy），$\rho_\pi(s_t, a_t)$ 为策略 $\pi$ 诱导的状态-动作边际分布，$\alpha > 0$ 为温度参数（Temperature Parameter），控制熵正则化的强度。将期望展开，目标函数等价地表示为：
\begin{equation}
    J_{\mathrm{MaxEnt}}(\pi) = \sum_{t=0}^{T} \mathbb{E}_{(s_t, a_t) \sim \rho_\pi} \left[ r(s_t, a_t) - \alpha \log \pi(a_t \mid s_t) \right]
\end{equation}

熵正则项的引入在理论上具有三重效应：其一，迫使策略在等效动作之间保持均匀分布，增强探索多样性；其二，使策略对奖励函数的微小扰动具有更强的鲁棒性；其三，为多峰奖励景观中的策略表示提供了自然的概率框架，有效防止策略过早收敛至次优确定性解。

\textbf{软贝尔曼方程}

在最大熵框架下，定义软状态价值函数（Soft State Value Function）$V^\pi_{\mathrm{soft}}$ 与软动作价值函数（Soft Q-Function）$Q^\pi_{\mathrm{soft}}$：
\begin{equation}
    V^\pi_{\mathrm{soft}}(s_t) = \mathbb{E}_{a \sim \pi(\cdot|s_t)} \left[ Q^\pi_{\mathrm{soft}}(s_t, a) - \alpha \log \pi(a \mid s_t) \right]
\end{equation}
\begin{equation}
    Q^\pi_{\mathrm{soft}}(s_t, a_t) = r(s_t, a_t) + \gamma \,\mathbb{E}_{s_{t+1} \sim \mathcal{P}} \!\left[ V^\pi_{\mathrm{soft}}(s_{t+1}) \right]
\end{equation}

将两式联立，得到软 Q 函数所满足的软贝尔曼期望方程（Soft Bellman Expectation Equation）：
\begin{equation}
    Q^\pi_{\mathrm{soft}}(s_t, a_t) = r(s_t, a_t) + \gamma \,\mathbb{E}_{\substack{s_{t+1} \sim \mathcal{P} \\ a_{t+1} \sim \pi}} \!\left[ Q^\pi_{\mathrm{soft}}(s_{t+1}, a_{t+1}) - \alpha \log \pi(a_{t+1} \mid s_{t+1}) \right]
\end{equation}

与标准贝尔曼方程相比，软贝尔曼方程在目标值中额外纳入了策略熵项，使得 Q 值在反映动作回报潜力的同时，亦对策略的随机性给予额外奖励。

\textbf{SAC 算法的参数化与损失函数}

Soft Actor-Critic（SAC）算法\cite{haarnoja2018soft}将最大熵强化学习框架与离策略（Off-Policy）深度学习相结合，通过维护两个 Critic 网络、一个 Actor 网络与对应的目标网络实现高效稳定的连续控制。

\underline{Critic 损失函数}：SAC 采用双 Q 网络（Double Q-Network）结构以缓解值函数高估问题。参数分别为 $\theta_1$ 和 $\theta_2$ 的两个 Critic 网络的训练目标定义为：
\begin{equation}
    L_Q(\theta_i) = \mathbb{E}_{(s_t, a_t, r_t, s_{t+1}) \sim \mathcal{D}} \left[ \left( Q_{\theta_i}(s_t, a_t) - y_t \right)^2 \right], \quad i \in \{1, 2\}
\end{equation}
其中，TD 目标 $y_t$ 由目标网络（参数分别为 $\bar{\theta}_1$，$\bar{\theta}_2$）与当前策略联合确定：
\begin{equation}
    y_t = r_t + \gamma \left( \min_{j \in \{1,2\}} Q_{\bar{\theta}_j}(s_{t+1}, \tilde{a}_{t+1}) - \alpha \log \pi_\phi(\tilde{a}_{t+1} \mid s_{t+1}) \right), \quad \tilde{a}_{t+1} \sim \pi_\phi(\cdot \mid s_{t+1})
\end{equation}

取两个目标网络 Q 值之最小值，有效抑制了函数逼近误差的累积传播，提升了 Critic 估计的保守性与训练稳定性。

\underline{Actor 损失函数}：为使 Actor 能够在连续动作空间中进行梯度优化，SAC 采用重参数化技巧（Reparameterization Trick）。将动作 $\tilde{a}$ 表示为噪声变量 $\varepsilon$ 的确定性函数 $\tilde{a}_t = f_\phi(\varepsilon_t;\, s_t)$，其中 $\varepsilon_t \sim \mathcal{N}(\mathbf{0}, \mathbf{I})$，$f_\phi$ 通常取压缩变换：
\begin{equation}
    \tilde{a}_t = \tanh\!\left( \mu_\phi(s_t) + \sigma_\phi(s_t) \odot \varepsilon_t \right)
\end{equation}

Actor 的训练目标为最小化以下损失：
\begin{equation}
    L_\pi(\phi) = \mathbb{E}_{s_t \sim \mathcal{D},\, \varepsilon_t \sim \mathcal{N}} \left[ \alpha \log \pi_\phi\!\left( f_\phi(\varepsilon_t;\, s_t) \mid s_t \right) - \min_{j \in \{1,2\}} Q_{\theta_j}\!\left( s_t,\, f_\phi(\varepsilon_t;\, s_t) \right) \right]
\end{equation}

通过重参数化，梯度可以直接从 Q 网络反向传播至 Actor 参数 $\phi$，将高方差的策略梯度估计转化为低方差的路径梯度估计，从而显著提升优化效率。

\underline{温度参数自适应调节}：SAC 将温度参数 $\alpha$ 视为可学习变量，通过最小化以下目标自动确定其值：
\begin{equation}
    L(\alpha) = \mathbb{E}_{a_t \sim \pi_\phi(\cdot|s_t)} \left[ -\alpha \log \pi_\phi(a_t \mid s_t) - \alpha \bar{\mathcal{H}} \right]
\end{equation}
其中，$\bar{\mathcal{H}}$ 为预设的目标熵，通常设为 $\bar{\mathcal{H}} = -\dim(\mathcal{A})$。当当前策略熵低于 $\bar{\mathcal{H}}$ 时，$\alpha$ 自动增大以加强探索；反之则减小以强化利用，实现了探索-利用权衡的动态自适应调节。

\textbf{探索与利用权衡的数学本质}

SAC 在探索-利用权衡（Exploration-Exploitation Trade-off）上的优越性，源自最大熵目标对最优策略解的结构性约束。可以证明，在最大熵框架下，最优软策略具有如下闭合形式：
\begin{equation}
    \pi^*_{\mathrm{soft}}(a \mid s) \propto \exp\!\left( \frac{1}{\alpha} Q^*_{\mathrm{soft}}(s, a) \right)
\end{equation}

此结论表明，最优软策略本质上是软 Q 值的玻尔兹曼分布（Boltzmann Distribution）。当 $\alpha \to 0$ 时，策略退化为确定性贪婪策略，对应纯利用；当 $\alpha \to \infty$ 时，策略趋于均匀分布，对应纯探索。温度参数 $\alpha$ 在两极端之间连续插值，提供了比 $\varepsilon$-贪婪等启发式方法更为严格的数学基础，并赋予探索行为以明确的信息论意义。

综上所述，SAC 通过最大熵目标、双 Q 网络估计与重参数化策略梯度的协同设计，在连续控制任务中实现了样本效率与训练稳定性的显著提升。然而，上述所有机制均建立在 Critic 和 Actor 网络能够从输入观测中提取高质量状态表征这一前提之上。在部分可观测场景下，若观测 $o_t$ 与真实状态 $s_t$ 之间存在较大的信息鸿沟，则值函数估计的精度与策略梯度的可靠性均将受到根本性制约。如何从高维、噪声观测中构建紧凑而充分的潜在状态表征，是提升上述算法在复杂视觉控制任务中性能的核心问题，这也正是本文后续章节所要深入探讨的主旨所在。

\xsection{离线强化学习与保守价值评估}{Offline Reinforcement Learning and Conservative Value Estimation}

在线强化学习（Online RL）要求智能体与环境持续交互以收集训练数据，然而在众多现实场景中——例如医疗决策、自动驾驶或机器人操作——此类交互的代价极为高昂甚至存在安全风险。离线强化学习范式的提出，旨在仅利用历史数据驱动策略学习，彻底规避在线交互的代价。然而，这一范式转换带来了标准强化学习框架所不具备的核心理论难题：分布偏移与外推误差。本节将系统阐述上述问题的数学本质，并深入剖析以隐式 Q 学习为代表的保守价值评估方法如何在理论层面予以解决。

\xsubsection{离线强化学习范式与分布偏移问题}{Offline Reinforcement Learning Paradigm and Distribution Shift}

\textbf{离线强化学习的基本范式}

离线强化学习（Offline Reinforcement Learning）的核心约束在于：学习过程中智能体无法与环境产生任何新的交互，策略优化所依赖的全部信息来源于一个预先收集的静态数据集：
\begin{equation}
    \mathcal{D} = \left\{ \left( s_t^{(i)},\, a_t^{(i)},\, r_t^{(i)},\, s_{t+1}^{(i)} \right) \right\}_{i=1}^{N}
\end{equation}
该数据集由一个或多个行为策略（Behavior Policy）$\pi_\beta$ 与环境交互产生，$\pi_\beta$ 可以是人类专家、启发式控制器、次优在线智能体或上述多者的混合。数据集中的状态-动作对服从行为策略诱导的分布 $\rho^{\mathcal{D}}(s, a) = d^{\pi_\beta}(s)\,\pi_\beta(a \mid s)$，其中 $d^{\pi_\beta}(s)$ 为 $\pi_\beta$ 的折扣状态访问分布。

在此约束下，离线强化学习的目标等价于：在 $\mathcal{D}$ 的支撑范围内寻找能够最大化真实环境回报的策略 $\pi$，而不依赖于额外的环境反馈。

\textbf{分布偏移的数学表征}

分布偏移（Distribution Shift）是离线强化学习区别于在线学习的根本性挑战。其核心矛盾在于：学习得到的策略 $\pi$ 在部署阶段所访问的状态-动作分布 $\rho^\pi(s, a) = d^\pi(s)\,\pi(a \mid s)$，与数据集所反映的分布 $\rho^{\mathcal{D}}(s, a)$ 之间存在结构性差异。

具体地，将两者之间的分布偏移量化为全变差（Total Variation）距离或 KL 散度：
\begin{equation}
    D_{\mathrm{TV}}\!\left( \rho^\pi \,\|\, \rho^{\mathcal{D}} \right) = \frac{1}{2} \int_{\mathcal{S} \times \mathcal{A}} \left| \rho^\pi(s, a) - \rho^{\mathcal{D}}(s, a) \right| \mathrm{d}s\,\mathrm{d}a
\end{equation}

当 $\pi$ 偏离 $\pi_\beta$ 时，$\rho^\pi$ 将在 $\mathcal{D}$ 未充分覆盖的区域赋予显著概率质量。正是在这些分布外（Out-of-Distribution, OOD）区域，Q 函数的泛化行为无法由数据监督，从而引发估计失真。

\textbf{外推误差的逻辑链及其灾难性后果}

将分布偏移导致值函数失效的过程形式化为以下推导链。

标准离线 Q 学习的贝尔曼更新为：
\begin{equation}
    Q_{\theta}(s, a) \leftarrow r(s, a) + \gamma \max_{a' \in \mathcal{A}} Q_{\bar{\theta}}(s', a')
\end{equation}
其中 $Q_{\bar{\theta}}$ 为目标网络。注意，上式中的 $\max_{a'}$ 算子在整个动作空间 $\mathcal{A}$ 上取最大，而非仅在数据集支撑 $\mathrm{supp}(\pi_\beta(\cdot \mid s'))$ 上取最大。

设 $a^*_{s'} = \arg\max_{a'} Q_{\bar{\theta}}(s', a')$ 为所选择的贪婪动作，当 $a^*_{s'} \notin \mathrm{supp}(\pi_\beta(\cdot \mid s'))$ 时，即 $a^*_{s'}$ 为分布外动作，神经网络函数逼近器对其 Q 值的估计完全依赖于在分布内数据上学到的归纳偏置（Inductive Bias）的外推，缺乏任何真实环境反馈的约束。

定义外推误差为：
\begin{equation}
    \varepsilon_{\mathrm{extrap}}(s, a) \triangleq \left| Q_\theta(s, a) - Q^{\pi_\beta}(s, a) \right|, \quad (s, a) \notin \mathrm{supp}(\mathcal{D})
\end{equation}

由于神经网络倾向于对分布外输入输出异常大的激活值，$Q_\theta$ 在 OOD 区域通常呈现严重的高估（Overestimation Bias）。更为致命的是，这一误差通过贝尔曼自举（Bellman Bootstrapping）的递归传播机制扩散至分布内的状态-动作对。设第 $k$ 次迭代后的误差上界为 $\varepsilon_k$，则有如下误差传播不等式：
\begin{equation}
    \varepsilon_{k+1} \leq \gamma \varepsilon_k + \gamma \cdot \varepsilon_{\mathrm{extrap}}
\end{equation}

该递推关系表明，外推误差以几何级数向历史状态-动作对反向传播，在迭代过程中造成系统性的 Q 值虚高。策略 $\pi$ 在此虚高 Q 值指引下倾向于选择 OOD 动作（因其 Q 值估计最大），而这些动作在真实环境中的回报往往远低于预期，最终导致训练完成后策略性能严重下降甚至彻底失效。

上述分析揭示，离线 Q 学习的失效根源并非算法设计的缺陷，而是函数逼近与分布偏移在贝尔曼自举框架下的结构性冲突。有效的解决方案须从根本上消除或约束对 OOD 动作 Q 值的依赖，这正是以 IQL 为代表的保守价值评估方法的理论出发点。

\xsubsection{隐式 Q 学习算法}{Implicit Q-Learning}

\textbf{理论初衷：消除分布外动作查询}

隐式 Q 学习（Implicit Q-Learning, IQL）\cite{kostrikov2021offline}的核心理论贡献在于，提出了一种完全规避对分布外动作进行显式 Q 值查询的价值评估范式。其关键洞察在于：贝尔曼最优算子中的 $\max_{a'} Q(s', a')$ 操作并非不可替代——若能直接估计数据集支撑范围内的最优状态价值函数 $V(s')$，则该项可以被安全地替换为 $V(s')$，从而在 TD 目标的计算中完全消除对具体动作的依赖。

然而，直接从数据集中估计 $V^*(s) = \max_{a \in \mathrm{supp}(\pi_\beta(\cdot|s))} Q^*(s, a)$ 面临如下困难：数据集中对于给定状态 $s$，动作的采样密度有限，简单取最大值将导致严重的有限样本偏差（Finite-Sample Bias）。IQL 通过期望回归（Expectile Regression）提供了这一问题的严格统计解。

\textbf{期望回归的数学推导}

设随机变量 $X$ 服从分布 $\mathcal{P}_X$，对于给定的期望水平 $\tau \in (0, 1)$，$X$ 的 $\tau$-期望值（$\tau$-Expectile）定义为以下优化问题的解：
\begin{equation}
    m_\tau = \mathop{\arg\min}_{m \in \mathbb{R}}\; \mathbb{E}_{X \sim \mathcal{P}_X}\!\left[ \mathcal{L}_2^\tau(X - m) \right]
\end{equation}
其中，非对称平方损失函数（Asymmetric Squared Loss）$\mathcal{L}_2^\tau: \mathbb{R} \to \mathbb{R}_{\geq 0}$ 定义为：
\begin{equation}
    \mathcal{L}_2^\tau(u) = \left| \tau - \mathbb{I}(u < 0) \right| u^2
\end{equation}
式中，$\mathbb{I}(\cdot)$ 为指示函数（Indicator Function）。展开可得：
\begin{equation}
    \mathcal{L}_2^\tau(u) = \begin{cases} \tau \, u^2, & u \geq 0 \\ (1 - \tau)\, u^2, & u < 0 \end{cases}
\end{equation}

对期望取一阶最优性条件，令 $\frac{\partial}{\partial m}\mathbb{E}[\mathcal{L}_2^\tau(X - m)] = 0$，得：
\begin{equation}
    2\tau \,\mathbb{E}\!\left[ (X - m)\,\mathbb{I}(X \geq m) \right] - 2(1-\tau)\,\mathbb{E}\!\left[ (X - m)\,\mathbb{I}(X < m) \right] = 0
\end{equation}

分析上式关于 $\tau$ 的极端情形：
\begin{itemize}
    \item 当 $\tau = 0.5$ 时，损失函数退化为标准均方误差，$m_{0.5}$ 为分布 $\mathcal{P}_X$ 的均值 $\mathbb{E}[X]$。
    \item 当 $\tau \to 1$ 时，正残差 $(X - m)$ 所受惩罚趋向无穷大，迫使 $m$ 趋近于 $X$ 的支撑上确界，即 $m_\tau \to \mathrm{ess\,sup}\,X$。
    \item 当 $\tau \in (0.5, 1)$ 时，$m_\tau$ 为均值与上确界之间关于 $\tau$ 单调递增的插值，即分布的高分位数估计。
\end{itemize}

这一性质表明，通过选取较大的 $\tau$ 值（如 $\tau = 0.7 \sim 0.9$），期望回归可以在统计意义上稳健地逼近分布的上端，从而隐式实现对策略价值上界的估计。

\textbf{IQL 的三阶段训练机制}

IQL 将价值估计解耦为三个独立的监督学习目标，各阶段均严格约束在数据集分布内。

\underline{第一阶段——状态价值函数的期望回归训练}：利用期望回归训练参数为 $\psi$ 的状态价值网络 $V_\psi(s)$，其损失函数为：
\begin{equation}
    L_V(\psi) = \mathbb{E}_{(s, a) \sim \mathcal{D}}\!\left[ \mathcal{L}_2^\tau\!\left( Q_{\bar{\theta}}(s, a) - V_\psi(s) \right) \right]
\end{equation}
其中，$Q_{\bar{\theta}}$ 为目标 Q 网络（参数 $\bar{\theta}$ 为 $\theta$ 的指数移动平均）。$V_\psi(s)$ 的物理意义得以精确刻画：它并非在整个动作空间上的期望或最大值，而是在给定状态 $s$ 下，数据集所记录的 Q 值分布 $\{Q_{\bar{\theta}}(s, a) : (s, a) \in \mathcal{D}\}$ 的 $\tau$-期望值——即数据集支撑范围内的高分位价值估计。由于 $V_\psi(s)$ 的优化目标仅涉及数据集中已有的 $(s, a)$ 对，整个过程完全不查询任何分布外动作的 Q 值。

\underline{第二阶段——Q 函数的保守贝尔曼更新}：利用第一阶段学得的 $V_\psi$ 替代标准贝尔曼算子中的 $\max_{a'} Q(s', a')$，对参数为 $\theta$ 的 Q 网络执行保守的 TD 更新：
\begin{equation}
    L_Q(\theta) = \mathbb{E}_{(s, a, r, s') \sim \mathcal{D}}\!\left[ \left( Q_\theta(s, a) - \underbrace{\left( r + \gamma V_\psi(s') \right)}_{\text{无动作查询的 TD 目标}} \right)^2 \right]
\end{equation}
TD 目标 $r + \gamma V_\psi(s')$ 在计算中无需指定动作 $a'$，完全规避了对 $s'$ 处任何具体动作 Q 值的查询，从根本上切断了外推误差的传播路径。

\underline{第三阶段——优势加权回归策略提取}：利用前两阶段的估计结果，通过优势加权回归（Advantage Weighted Regression, AWR）\cite{peng2019advantage}从数据集中提取策略 $\pi_\phi$：
\begin{equation}
    L_\pi(\phi) = -\mathbb{E}_{(s, a) \sim \mathcal{D}}\!\left[ \exp\!\left( \beta \cdot A(s, a) \right) \cdot \log \pi_\phi(a \mid s) \right]
\end{equation}
其中，优势函数 $A(s, a) = Q_{\bar{\theta}}(s, a) - V_\psi(s)$ 衡量数据集中动作 $a$ 相对于当前 $V_\psi$ 估计值的超额回报，$\beta > 0$ 为逆温度参数。

AWR 目标的等价形式为加权最大似然估计：
\begin{equation}
    L_\pi(\phi) = -\mathbb{E}_{(s, a) \sim \mathcal{D}}\!\left[ w(s, a) \cdot \log \pi_\phi(a \mid s) \right], \quad w(s, a) = \exp\!\left( \beta \bigl( Q_{\bar{\theta}}(s, a) - V_\psi(s) \bigr) \right)
\end{equation}

权重 $w(s, a)$ 对优势为正的动作（即优于当前值函数估计的动作）赋予指数级放大的模仿权重，对优势为负的动作赋予接近零的权重，从而在纯行为克隆（Behavioral Cloning）与纯策略优化之间取得原则性权衡。通过参数 $\beta$ 的调节，IQL 可以在"保守地模仿数据集"与"激进地提取数据集最优行为"之间灵活权衡，而整个提取过程始终限制在数据集的动作分布支撑范围内，有效规避了分布外泛化的风险。

\textbf{IQL 的理论一致性}

IQL 的三阶段设计具有严格的理论一致性保证。当数据集覆盖充分（$\mathcal{D}$ 包含最优策略的支撑）且 $\tau \to 1$ 时，$V_\psi(s)$ 在统计意义上趋近于 $\max_{a \in \mathrm{supp}(\pi_\beta(\cdot|s))} Q^*(s, a)$，Q 网络的 TD 更新趋近于贝尔曼最优算子，AWR 策略提取趋近于最优策略。在 $\tau < 1$ 的一般情形下，IQL 提供了一个关于 $\tau$ 连续参数化的保守性-最优性权衡谱系，为实践中针对数据集质量的自适应调节提供了理论依据。

综上所述，IQL 通过将最优价值评估问题转化为期望回归——一种严格约束于数据集分布的统计估计任务——从根本上消除了离线 Q 学习对分布外动作查询的依赖，并借助优势加权回归实现了稳健而高效的策略提取。这一价值评估范式的转变表明：在复杂观测场景下，高质量的状态表征对于准确估计 $Q_{\bar{\theta}}(s, a)$ 与 $V_\psi(s)$ 同等关键——若状态编码器无法从高维噪声观测中提取充分的语义特征，则期望回归所估计的 $\tau$-期望值将因输入信息不充分而产生系统性偏差，进而损害策略提取的质量。这一理论分析为后文在复杂残缺观测环境下进一步优化特征表征的研究路径提供了明确的方法论基础。

\xsection{复杂观测下的表示学习机制}{Representation Learning Mechanisms under Complex Observations}

前两节分别从决策框架与价值评估角度建立了强化学习的理论基础。然而，上述框架均隐含一个前提：智能体能够从原始观测中获取足以支撑精确价值估计的状态信息。当观测为高维图像或存在结构性缺失时，这一前提不再自然成立，表征学习的质量成为制约整体性能的关键瓶颈。本节从两个互补维度剖析复杂观测场景下的表征构建机制：其一，基于数据增强的视觉不变性表征；其二，基于掩码建模的低维物理状态重构。两者共同构成本文核心算法设计的表征理论基础。

\xsubsection{视觉强化学习与数据增强表示}{Visual Reinforcement Learning and Data Augmentation Representation}

\textbf{卷积神经网络的降维表征作用}

基于图像的视觉强化学习（Visual Reinforcement Learning）面临的首要挑战是维度灾难（Curse of Dimensionality）：原始 RGB 观测 $o_t \in \mathbb{R}^{H \times W \times C}$ 的像素空间维度通常达到 $10^4 \sim 10^5$ 量级，而真正决定系统动力学的物理自由度（关节角、速度、接触力等）往往仅有数十维。卷积神经网络（Convolutional Neural Network, CNN）凭借局部感受野与参数共享的归纳偏置，对视觉观测实施层次化的非线性特征提取：
\begin{equation}
    z_t = E_\phi(o_t), \quad E_\phi: \mathbb{R}^{H \times W \times C} \to \mathbb{R}^d,\quad d \ll H \cdot W \cdot C
\end{equation}
其中，$E_\phi$ 为参数化的 CNN 编码器，$z_t \in \mathbb{R}^d$ 为低维潜特征（Latent Feature）。CNN 通过卷积层逐步聚合局部纹理、边缘与形状信息，经全连接层投影至紧凑的潜在空间，从而将高维像素观测压缩为适合值函数与策略网络处理的低维表征。然而，潜特征 $z_t$ 的语义质量直接决定了下游强化学习算法的性能上限，这催生了对编码器训练目标的精细设计。

\textbf{不变性表征的数学定义}

数据增强（Data Augmentation）是提升视觉编码器鲁棒性的主流手段。其理论动机来源于不变性表征（Invariant Representation）的概念：对于任务无关的观测变换（如随机裁剪、颜色抖动、平移或频率掩盖），编码器所提取的潜特征应在语义层面保持稳定，不受变换引起的像素级扰动影响。

设 $\mathcal{T}$ 为一族数据增强变换的分布，$\tilde{o} = \mathcal{T}(o)$ 为经增强操作后的观测。不变性表征的目标可形式化表述为：
\begin{equation}
    \forall \mathcal{T} \in \mathbb{T}:\quad \mathcal{F}\!\left( E_\phi(o) \right) \approx \mathcal{F}\!\left( E_\phi\!\left( \mathcal{T}(o) \right) \right)
\end{equation}
其中，$\mathcal{F}$ 为任意下游任务相关的泛函（如价值函数或策略映射）。在实践中，不变性目标通常通过对比学习（Contrastive Learning）或一致性正则化（Consistency Regularization）来实现，例如 RAD\cite{laskin2020reinforcement}、DrQ\cite{kostrikov2020image} 等方法要求同一观测的两个增强视图在潜空间中距离足够近：
\begin{equation}
    \mathcal{L}_{\mathrm{inv}} = \mathbb{E}_{o \sim \mathcal{D},\, \mathcal{T}_1, \mathcal{T}_2 \sim \mathbb{T}}\!\left[ \left\| E_\phi\!\left( \mathcal{T}_1(o) \right) - \mathrm{sg}\!\left[ E_\phi\!\left( \mathcal{T}_2(o) \right) \right] \right\|_2^2 \right]
\end{equation}
其中，$\mathrm{sg}[\cdot]$ 表示停止梯度算子（Stop-Gradient），用于防止表征坍缩（Representation Collapse）。不变性约束迫使编码器剥离增强操作引入的无关扰动，保留对决策任务真正有用的结构信息。

\textbf{关键痛点：强数据增强对时间平滑性的破坏}

尽管数据增强能够改善编码器对视觉扰动的鲁棒性，强数据增强（Strong Augmentation）——即变换幅度显著改变图像的低频结构，如大范围随机裁剪、频率域掩盖或强烈颜色变换——在强化学习中会引发一个深层的理论问题：马尔可夫决策过程的时间平滑性（Temporal Smoothness）遭受破坏。

时间平滑性要求相邻时间步的状态转移在潜空间中保持局部连续性，即对于真实状态序列 $(s_t, s_{t+1})$，其编码表征满足：
\begin{equation}
    \left\| E_\phi(o_{t+1}) - E_\phi(o_t) \right\|_2 \leq \kappa \left\| o_{t+1} - o_t \right\|_2
\end{equation}
其中，$\kappa > 0$ 为利普希茨常数（Lipschitz Constant），刻画了编码器的局部平滑程度。

当施加强数据增强 $\mathcal{T}$ 时，同一时间步的观测 $o_t$ 被映射至两个语义相异的增强视图 $\tilde{o}_t^{(1)}$ 与 $\tilde{o}_t^{(2)}$，而相邻时间步 $\tilde{o}_t$ 与 $\tilde{o}_{t+1}$ 因独立采样增强操作而彼此无关。在此情形下，时序差分（Temporal Difference）自举所依赖的目标值为：
\begin{equation}
    y_t = r_t + \gamma \min_{j} Q_{\bar{\theta}_j}\!\left( E_\phi\!\left( \mathcal{T}^{(j)}(o_{t+1}) \right),\, \tilde{a}_{t+1} \right)
\end{equation}
其中 $\mathcal{T}^{(j)}$ 为对 $o_{t+1}$ 的第 $j$ 次随机采样。目标值 $y_t$ 的方差可分解为：
\begin{equation}
    \mathrm{Var}(y_t) = \gamma^2\, \mathrm{Var}\!\left( Q_{\bar{\theta}}\!\left( E_\phi(\mathcal{T}(o_{t+1})), \tilde{a}_{t+1} \right) \right) \geq \gamma^2\, \mathbb{E}\!\left[ \left( \frac{\partial Q}{\partial z} \right)^2 \right] \cdot \mathrm{Var}\!\left( E_\phi\!\left(\mathcal{T}(o_{t+1})\right) \right)
\end{equation}

上式表明，强增强所引入的编码方差 $\mathrm{Var}(E_\phi(\mathcal{T}(o_{t+1})))$ 经 Q 函数梯度的放大后，直接转化为目标值 $y_t$ 的高方差。目标值的随机波动等价于在贝尔曼更新中注入系统性噪声，致使 Critic 的梯度方向在相邻迭代间剧烈震荡，训练曲线呈现高度不稳定性，甚至导致 Q 值发散。

此问题的根源在于：数据增强操作集合 $\mathbb{T}$ 的设计与具体强化学习任务的奖励结构、状态转移动力学之间存在天然的语义错位——增强对于分类任务可能是无害的标签保持变换，但对于依赖精确时序信息进行 TD 自举的强化学习任务，相同的增强却可能破坏决策所必需的物理一致性。这一矛盾揭示了一个核心需求：针对强化学习任务，数据增强的选择不应遵循与视觉识别任务相同的启发式规则，而需要一种能够感知当前任务动力学属性的增强选择机制，以在不变性收益与时间平滑性代价之间取得原则性权衡。这一需求正是本文后续提出 PICK 算法的直接理论动因。

\xsubsection{掩码建模与自监督特征重建}{Masked Modeling and Self-Supervised Feature Reconstruction}

\textbf{数据流形假设}

高维视觉观测的可学习性根植于数据流形假设（Data Manifold Hypothesis）：自然场景中由物理系统产生的图像序列，尽管在像素空间中具有极高维度，却实际分布在一个远低维的黎曼流形（Riemannian Manifold）$\mathcal{M} \subset \mathbb{R}^{H \times W \times C}$ 附近，且该流形的内在维度（Intrinsic Dimensionality）由物理系统的真实自由度数目所决定：
\begin{equation}
    \dim(\mathcal{M}) \approx \dim(\mathcal{S}) \ll H \cdot W \cdot C
\end{equation}

这一假设的物理内涵在于：机器人系统的关节角度、速度、接触力等状态变量之间受到运动学约束、动力学方程与接触力学定律的严格约束，使得合法的状态轨迹被限制在物理上可行的低维子流形上。流形上的相邻点对应物理上连续的状态演化，流形的切空间局部描述了合法的状态变化方向。在此假设下，表征学习的目标等价于：寻找从像素空间到流形内在坐标的高保真映射，从而在低维潜空间中保留物理系统的几何结构与动力学约束。

\textbf{掩码自编码器的架构与重建目标}

掩码自编码器（Masked Autoencoder, MAE）\cite{he2022masked}是实现上述流形学习目标的有效自监督框架。其核心思想源于完形填空（Cloze Test）的认知原理：若模型能够从部分遮盖的输入中精确重建被掩盖的区域，则必然已学得数据的全局结构与局部语义之间的内在关联，这恰好对应于对流形几何的隐式编码。

MAE 由编码器 $E_\phi: \mathbb{R}^{H \times W \times C} \to \mathbb{R}^d$ 与解码器 $D_\psi: \mathbb{R}^d \to \mathbb{R}^{H \times W \times C}$ 构成。给定原始观测 $s \in \mathbb{R}^{H \times W \times C}$，首先施加随机掩码操作 $M \in \{0, 1\}^{H \times W}$（$M_{ij} = 0$ 表示该位置被遮盖），得到部分可见输入：
\begin{equation}
    \tilde{s} = s \odot M
\end{equation}
其中，$\odot$ 表示逐元素乘积（Hadamard Product），实践中 $M$ 通常对图像块（Patch）级别进行采样，掩码比例 $\rho \in [0.5, 0.75]$ 以确保任务难度足够。

编码器对可见区域进行特征提取，解码器以潜特征与可学习掩码标记（Mask Token）为输入，尝试重建完整观测。训练目标为最小化重建均方误差：
\begin{equation}
    \min_{\phi,\, \psi}\;\mathbb{E}_{s \sim \mathcal{D},\, M \sim \mathcal{M}}\!\left[ \left\| D_\psi\!\left( E_\phi(s \odot M) \right) - s \right\|_2^2 \right]
\end{equation}

在实践中，损失计算通常仅限于被掩盖区域以提高学习效率：
\begin{equation}
    \mathcal{L}_{\mathrm{recon}}(\phi, \psi) = \mathbb{E}_{s \sim \mathcal{D},\, M \sim \mathcal{M}}\!\left[ \frac{1}{|\bar{M}|}\sum_{(i,j):\, M_{ij}=0} \left\| \left[ D_\psi\!\left( E_\phi(s \odot M) \right) \right]_{ij} - s_{ij} \right\|_2^2 \right]
\end{equation}
其中，$|\bar{M}| = \sum_{ij}(1 - M_{ij})$ 为被掩盖区域的总面积。

\textbf{自监督重建的理论可行性}

MAE 重建目标的理论可行性建立在以下因果链上：其一，被掩盖区域的精确重建要求模型对未掩盖区域的特征进行"语境推理"（Contextual Reasoning），即理解相邻图像块之间的视觉依赖关系；其二，视觉图像块之间的依赖关系在物理系统中由系统的几何结构、光照模型与运动学约束决定，因此语境推理过程隐式地要求模型学习这些物理约束；其三，编码器 $E_\phi$ 为了支持解码器 $D_\psi$ 完成高质量重建，必须将可见区域的信息压缩为包含全局物理状态信息的紧凑表征 $z = E_\phi(\tilde{s})$，而非仅编码局部纹理特征。

这一自监督过程等价于将潜特征 $z$ 约束在数据流形 $\mathcal{M}$ 的内在坐标上：若 $z$ 偏离流形，则解码器 $D_\psi$ 无法从该坐标重建出物理合法的观测，训练损失将产生显著惩罚，迫使 $z$ 回归流形。形式上，最优编码器 $E_\phi^*$ 满足：
\begin{equation}
    E_\phi^*(s) \approx \Phi(s), \quad \forall s \in \mathcal{M}
\end{equation}
其中，$\Phi: \mathcal{M} \to \mathbb{R}^d$ 为流形 $\mathcal{M}$ 到潜空间的同胚映射（Homeomorphism），即一一连续的双向映射，保留了流形的拓扑结构与局部度量关系。因此，MAE 以纯数据驱动的方式，隐式学习了物理系统状态空间的几何编码，为下游强化学习中的价值函数估计提供了信息充分的低维特征基础。

\textbf{关键痛点：联合训练中的梯度冲突}

在强化学习框架下，一种直觉性的方法是将重建损失 $\mathcal{L}_{\mathrm{recon}}$ 作为辅助任务与 TD 损失 $\mathcal{L}_{\mathrm{TD}}$ 联合优化，对共享编码器 $E_\phi$ 进行端到端训练：
\begin{equation}
    \mathcal{L}_{\mathrm{joint}}(\phi) = \mathcal{L}_{\mathrm{TD}}(\phi) + \lambda \,\mathcal{L}_{\mathrm{recon}}(\phi, \psi)
\end{equation}

然而，这一联合训练范式在实践中面临严重的梯度冲突（Gradient Conflict）问题。设 $g_{\mathrm{TD}} = \nabla_\phi \mathcal{L}_{\mathrm{TD}}$ 与 $g_{\mathrm{recon}} = \nabla_\phi \mathcal{L}_{\mathrm{recon}}$ 分别为两个任务施加在编码器参数上的梯度，梯度冲突的严重程度由其余弦相似度刻画：
\begin{equation}
    \cos\angle(g_{\mathrm{TD}},\, g_{\mathrm{recon}}) = \frac{\langle g_{\mathrm{TD}},\, g_{\mathrm{recon}} \rangle}{\left\| g_{\mathrm{TD}} \right\|_2 \cdot \left\| g_{\mathrm{recon}} \right\|_2}
\end{equation}

当 $\cos\angle(g_{\mathrm{TD}}, g_{\mathrm{recon}}) < 0$ 时，两梯度方向相反，合并梯度 $g_{\mathrm{TD}} + \lambda g_{\mathrm{recon}}$ 既无法充分减小 $\mathcal{L}_{\mathrm{TD}}$ 也无法充分减小 $\mathcal{L}_{\mathrm{recon}}$，编码器陷入两个目标之间的博弈均衡，表征质量显著退化。

这一冲突的产生有其深刻的几何根源。TD 损失的梯度 $g_{\mathrm{TD}}$ 所传递的信号来自贝尔曼残差，其方向由当前 Q 值的高估或低估方向决定，本质上是一种间接的、通过奖励信号反向推断出的特征选择压力——它倾向于强化与当前奖励结构相关的特征维度，对价值估计无关的特征施加抑制。相比之下，重建损失的梯度 $g_{\mathrm{recon}}$ 来自像素级重建误差，所传递的是对全局物理状态信息的保留压力——它倾向于均匀地保留所有具有重建价值的视觉特征，包括背景纹理、光照变化等对决策无关但对像素重建有贡献的信息。

两者的目标在特征选择上存在本质矛盾：TD 学习驱使表征向任务特定的低维决策子空间坍缩，而重建学习驱使表征向全局物理状态的高维流形展开。若在共享编码器上直接叠加两种梯度，则重建目标提取的流形结构将被 TD 更新的噪声梯度所污染，编码器既无法学得用于精确值函数估计的任务相关特征，也无法维护用于重建的全局状态表征，最终导致两项任务的性能均低于各自独立训练的基线。

这一梯度冲突分析表明，自监督重建与 TD 优化在同一编码器参数空间的共存需要精细的梯度流控制机制。必须设计一种表征学习与决策优化解耦的训练范式——允许重建任务在编码器中自由构建物理状态的流形表征，同时防止 TD 梯度对该流形结构造成破坏性干扰。这一原则性要求正是本文后续提出 Masked-IQL 梯度截断机制的核心理论动因。
